% Review template preamble file
% Based on Paolo Adajar's template: https://github.com/padajar/paolo-pset/tree/main

%%%%%%%%%%
% COURSE-SPECIFIC INFORMATION
%%%%%%%%%%

% Edit these as-needed: these commands can be included in any TeX file in the directory

\newcommand{\name}{Isaac Baumann}
\newcommand{\email}{dbaum@uic.edu}
\newcommand{\classnum}{ECON 101}
\newcommand{\subject}{Principles of Microeconomics}
\newcommand{\instructors}{Dr. Johnson}
\newcommand{\semester}{Fall 2024}

%%%%%%%%%%
% PREAMBLE
%%%%%%%%%%

% Font setting
\usepackage{sansmathfonts}
\usepackage[T1]{fontenc}
\renewcommand*\familydefault{\sfdefault}

% Primary packages
\usepackage{lipsum}

\usepackage{enumitem}
\usepackage{amsmath}
\usepackage[dvipsnames]{xcolor}
\usepackage{amssymb}
\usepackage{fancyhdr}
\usepackage{tcolorbox}
\usepackage{hyperref}
\hypersetup{
    colorlinks,
    citecolor=black,
    filecolor=blue,
    linkcolor=black,
    urlcolor=blue
}
\usepackage{siunitx}

% Secondary packages: add/remove as needed
\usepackage{tikz}
    \usetikzlibrary{calc, patterns, angles, quotes, arrows.meta}
\usepackage{pgfplots}
\usepackage{multicol}				% Use multiple columns
    \columnsep=20pt				% Set space between columns in points
\usepackage[linesnumbered,ruled,vlined]{algorithm2e}
\usepackage[toc,page]{appendix}
\usepackage{siunitx}

% Set header formatting
\pagestyle{fancy}
\fancyhf{}
\renewcommand{\sectionmark}[1]{\markboth{#1}{}}
\fancyhead[L]{\leftmark}
\fancyhead[R]{\classnum}
\cfoot{\thepage}

% Custom functions
\newcommand\numberthis{\addtocounter{equation}{1}\tag{\theequation}}
\let\vec\mathbf
% \newcommand{\rank}{\text{Rank}}
\newcommand{\corr}{\text{Corr}}
% \newcommand{\tr}{\text{tr}}
\newcommand{\plimtxt}{\text{plim }}
\newcommand{\plim}{\xrightarrow{\text{p}}}
\newcommand{\dconv}{\xrightarrow{\text{d}}}
\newcommand{\exptfn}[1]{\mathbb{E}\left[{#1}\right]}
\newcommand{\varfn}[1]{\text{Var}\left({#1}\right)}
\newcommand{\covfn}[1]{\text{Cov}\left({#1}\right)}
\newcommand{\indep}{\perp \!\!\! \perp}

% Comment block
\newcommand{\commentblock}[1]{}